\section*{Summary}

At fixed coupling the rate of a crystal adds sector by sector, and which sector dominates follows from two facts: how much energy each sector can absorb, and whether the coupling is screened (\cref{sec:17.1}). A single phonon carries at most the top of the band, \SI{64.1}{meV} in silicon; a freely recoiling nucleus at most $E_R^{\mathrm{max}}$, \SI{0.054}{eV} at \SI{10}{MeV}; an electronic excitation up to the whole kinetic energy, \SI{35}{eV} at \SI{10}{MeV}. For a nucleon coupling with a heavy mediator the electron-following sector holds every threshold below the band top, with one phonon leading up to \SI{20}{MeV} at \SI{1}{meV}, and the Migdal effect holds the band between the end of the lattice spectrum and the largest deposit. The multiphonon--Migdal crossover threshold in silicon is 0.178, 0.307 and \SI{5.33}{eV} at 1, 10 and \SI{100}{MeV}: pinned near three times the band top for soft kicks, climbing through the Poisson tail of the multiphonon spectrum between 3 and \SI{20}{MeV}, and sitting on the free-recoil endpoint above the impulse onset. For an electron coupling, screening decides the low thresholds. A dark photon couples to the total charge, its structure factor is the loss function, and silicon, with vanishing Born charges, has no lattice channel for it at order $q^2$; a mediator coupled to the electrons alone sees the electrons that follow the nuclei and raises the rate at \SI{1}{MeV} from 17 to $7\times10^5$ events per kilogram and year at $\sigmae = \SI{e-37}{cm^2}$. In GaAs the Born charge $Z^* = 2.11$ reverses the picture: the polar phonon opens to the dark photon at \SI{10.35}{keV} and is nearly closed to a nucleon coupling.

For absorption each boson measures one response at $q = p_\phi$ and $\omega = m_\phi$ (\cref{sec:17.2}). The dark photon sees the loss function, which screening caps at $L \le 1/(2|\varepsilon_1|)$, so it is absorbed strongly near the zeros of $\varepsilon_1$: the longitudinal phonon of GaAs, the highest polar branch of sapphire and the plasmon of silicon. The axion-like particle sees the conductivity or the spin current and has no such ceiling; in a clean superconductor its spin current breaks pairs directly. Between 0.38 and \SI{1.17}{eV} no target of the book covers the dark photon at leading order. A counter of quanta whose band is the line beats a quantum-limited linear readout of the same mode by $\ee^{m_\phi/2T}$ in signal, so its advantage vanishes for masses well below the temperature, \SI{0.86}{\micro eV} at \SI{10}{mK}, and above twenty times the temperature the counter is limited by the absorbed quanta alone.

Dark matter with a mass splitting can lose its excitation energy to a collective mode (\cref{sec:17.3}). When the mode's group velocity exceeds the halo speeds and its weight is smooth over $\mchi v_{\mathrm{max}}$, the rate is independent of the halo to order $v^2/v_\nu^2$, the line sits where the mode's dispersion meets the shell of fixed momentum, and its width has the shape of $\eta$. In silicon at \SI{1}{MeV} a splitting of \SI{10}{eV} raises the single-phonon rate 31 times over elastic scattering, and the shell form gives 0.78 to 1.10 of the full integral. For GaAs at \SI{30}{keV} and $\Delta_\chi = -\SI{1}{eV}$, down-scattering into the longitudinal phonon gives 1810 events per kilogram and year against 842 for elastic scattering. An endothermic splitting of \SI{1}{eV} brings every collective channel to \SI{0.3}{MeV}, half the pair threshold of silicon, and the plasmon is closed at every mass.

Cross terms between sectors appear for couplings to both electrons and nuclei (\cref{sec:17.4}). The interference between the Migdal and the phonon-assisted amplitudes is bounded by the positivity of the response matrix, \cref{eq:17.4-bound}, and at small momentum transfer the bound is reached: the summed Migdal amplitude of a one-phonon pair state is $-(m_e/M_\kappa)$ times the phonon-assisted one, so the phonon-dressed rate carries the factor $|\tilde V_e - (m_e/M_\kappa)\tilde V_{\mathrm{nuc}}|^2$, \cref{eq:17.4-rate}. For a dark photon this raises the phonon-assisted rate of silicon by $5.47\times10^{-4}$; for a coupling in proportion to mass the bracket vanishes. In the windows of silicon and GaAs, at a fifth to three fifths of the way to $2\pi/a$, the relation fixes the cross term in order of magnitude.

Each channel then reduces to a few material numbers, each valid in a window (\cref{sec:17.5}): the floor $2\Eg/v_{\mathrm{max}}^2$, the inverse effective sound speed $1/c_{\mathrm{NR1}}$ for a heavy mediator, the acoustic edge, the polar band integral $\mathcal{I}_{\mathrm{pol}}/\rho_T$ and the polar edge, and the Migdal figure of merit. The reach figures of merit are kinematic and exact; the rate figures of merit hold in windows. No scalar ranks electron-recoil rates, and the partial $f$-sum of the valence electrons caps them through $\omega_p^{\mathrm{val}}$, \SI{16.6}{eV} in silicon. The Migdal figure of merit lies between a plasmon-pole limit and a limit with the weight at the gap, which differ by $5.7\times10^5$ in silicon, and the pole limit itself depends on the dispersion of the plasmon, so its ranking is open. A heavy-mediator ranking holds pairwise, in the window $W_{AB}$ of two crystals, because the shift of $\eta$ by the sound speed breaks a common window; at \SI{0.3}{MeV} the computed rate ratios of germanium and diamond to silicon are 1.73 and 0.46 against the speed ratios 1.69 and 0.49 (\cref{sec:17.5}).

Rankings by these numbers carry the uncertainty of their inputs (\cref{sec:17.6}). Each input gets a provenance width by how it was obtained, measured, computed or estimated; the widths are conventions, and the database sound speeds of silicon, germanium and diamond, low by 1 to \SI{10}{\percent} and all in the same direction, show both the size of one convention and why method errors correlate. Monte Carlo propagation gives the probability that a crystal ranks first and its interval of ranks; heights in the interval order of the 68\% intervals give tiers, which never assert an order the intervals do not contain. Among eleven crystals germanium leads the electron-recoil floor, TeO$_2$ the heavy-mediator single-phonon rate (probability 0.968, on a computed tensor), diamond the nuclear edge, LiF the polar rate and GaAs the polar reach. Silicon leads none. Stiff lattices and light polar ions win one row of each pair and lose the other. The Migdal channel has no reading at independent errors, its largest probability of ranking first being 0.315, and the reading it acquires when the pole prior is shared by all crystals is an artifact of that assumption. The light-mediator single-phonon rate does not rank crystals at all. A shared method error sharpens comparisons inside one method and cannot help a comparison of a measured with a computed value, which keeps the method's offset in full.

\Cref{tab:17-nuisances} collects the truncations and inputs of the chapter with their sizes. How the energy per quantum, the readout and the backgrounds change these rankings is the subject of the guide to next targets of \cref{sec:21.7}.

\begin{table}[tbp]
  \caption{Nuisances of the comparison across channels. Type in the classes of \cref{sec:8.8}: normalization, tilt, threshold and new support; ``order'' marks an item that acts on a ranking rather than on a rate. Basis of the size: ``computed'' (next order or exact result computed), ``model'' (computed in a model), ``est.''\ (parametric estimate), ``input spread'' (spread or limitation of a measured or computed input), ``convention'' (a width assigned by convention), ``unsized'' (no size known). The last column gives the section where the item is treated.}
  \label{tab:17-nuisances}
  \centering\scriptsize\setlength{\tabcolsep}{3pt}
  \begin{tabular}{@{}>{\raggedright\arraybackslash}p{0.26\linewidth}>{\raggedright\arraybackslash}p{0.12\linewidth}>{\raggedright\arraybackslash}p{0.35\linewidth}>{\raggedright\arraybackslash}p{0.09\linewidth}>{\raggedright\arraybackslash}p{0.10\linewidth}@{}}
    \toprule
    Truncation or input & Type & Size & Basis & Section \\
    \midrule
    Coherent corrections to the multiphonon tail & threshold & a factor of two in the tail moves the crossover in silicon at \SI{10}{MeV} by \SI{0.014}{eV} & model & \cref{sec:17.1} \\
    Local fields and off-diagonal elements of the Migdal probability & normali\-zation & diagonal factor 1.35 moves the crossover at \SI{10}{MeV} by \SI{0.006}{eV}; off-diagonal size open (\cref{sec:12.3}) & computed; unsized & \cref{sec:17.1} \\
    Coherent Migdal response at $q < 2q_{\mathrm{BZ}}$ & normali\-zation; tilt in $\omega_e$ & diagonal form $14.9\times$ high as $q \to 0$ (\cref{eq:8.3-boostward}), quantitative only below about \SI{0.75}{keV}; \SI{0.75}{keV} to $2q_{\mathrm{BZ}}$ unsized, nearly all of the silicon Migdal rate up to \SI{1}{MeV} (\SI{29}{\percent} of it from $q < q_{\mathrm{BZ}}$ at \SI{1}{MeV}), \SI{33}{\percent} at \SI{2}{MeV}, \SI{7}{\percent} at \SI{3}{MeV} (\cref{tab:12-nuisances}) & computed (limit); unsized & \cref{sec:17.1} \\
    Electron-following sector of an electron-only coupling (long-wavelength model) & normali\-zation; tilt in $\mchi$ & equality with electron recoil at \SI{12.3}{MeV} is a property of the model; Umklapp, transverse and multiphonon terms add, dispersion and loss of following charge subtract & unsized & \cref{sec:17.1} \\
    Phonon-assisted side of the silicon window & normali\-zation & estimate known to a factor of ten; lead to about \SI{0.43}{MeV}, two-channel band 0.43--\SI{0.51}{MeV} & est. & \cref{sec:17.1} \\
    Germanium zero-phonon edge and phonon density of states & threshold & not known; the germanium electronic window and crossover are not drawn & unsized & \cref{sec:17.1} \\
    Half-gap and cutoff of ZrTe$_5$ & threshold & measured-structure (\SI{11.8}{meV}) and relaxed-cell (\SI{15.6}{meV}) parameters move the absorption onset from 23.6 to \SI{31.2}{meV}; the reach at \SI{100}{meV} changes by \SI{2.5}{\percent} & input spread & \cref{sec:17.2} \\
    Counting against linear readout & normali\-zation & a counter band $B$ wider than the line multiplies the advantage by $(B/\Delta\omega_{\mathrm{eq}})^{-1/2}$; squeezed readout and dark counts reduce it & unsized & \cref{sec:17.2} \\
    Exothermic shell form & normali\-zation & halo-independent to $O(v^2/v_\nu^2)$; silicon phonons 0.78--1.10 of the full integral at \SI{1}{MeV}, the weight varying on the zone scale & computed & \cref{sec:17.3} \\
    Electron-gas plasmon against crystal tables & normali\-zation & plasmon-pole rate 0.09--0.16 of the tables' electronic rate at 18--\SI{22}{eV} (two models) & model & \cref{sec:17.3} \\
    Migdal--phonon-assisted cross term & normali\-zation & bounded by positivity, \cref{eq:17.4-bound}; exact at small $q$ (\cref{eq:17.4-rate}: $+5.47\times10^{-4}$ for a dark photon in silicon); in the windows, at 0.2--0.6 of $2\pi/a$, order of magnitude only & computed (limit); est. & \cref{sec:17.4} \\
    Shift of $\eta$ by the sound speed (heavy mediator) & tilt in $\mchi$; order & rate ratios differ from the speed ratios by up to \SI{7}{\percent} at \SI{0.3}{MeV} (germanium, diamond against silicon); rankings hold pairwise in $W_{AB}$ & computed & \cref{sec:17.5} \\
    Interband weight of the Migdal pole factor & normali\-zation; order & pole and gap limits differ by $5.7\times10^5$ in silicon; the dispersion of the pole changes LiF over silicon from 0.87 to 1.4--1.7; the ranking uses the pole form with a prior $X_{\mathrm{pole}}$ log-uniform on [1, 3], a device (silicon's own value is about 25), and gives no reading at independent errors & est. & \cref{sec:17.5,sec:17.6} \\
    Provenance widths & order & database speeds 0.086 (three checked crystals: root mean square 0.065); gap shift 0.03--\SI{0.10}{eV}; doubling the speed width lowers TeO$_2$'s probability of ranking first from 0.968 to 0.781 & convention & \cref{sec:17.6} \\
    Offset of the computed elastic tensors & order & database speeds low by \SI{4.9}{\percent} on average (silicon, germanium, diamond); GaAs against germanium 0.887 $\to$ 0.744 with the offset & input spread & \cref{sec:17.6} \\
    Method correlation $r_{\mathrm{mat}}$ & order & sharpens one-method comparisons (TeO$_2$ against GaAs 0.976 at 0, 1.000 at 1); resolves the Migdal tie artificially (NaF 0.315 $\to$ 1.000) & computed & \cref{sec:17.6} \\
    Room-temperature gaps used at low temperature & threshold; order & germanium against GaSb 0.973 with the shift 0.03--\SI{0.10}{eV}, 0.837 with 0--\SI{0.15}{eV}; reversed without shift & est. & \cref{sec:17.6} \\
    \bottomrule
  \end{tabular}
\end{table}
